\documentclass[11pt,a4paper]{article}

\usepackage[T1]{fontenc}
\usepackage[utf8]{inputenc}
\usepackage{lmodern}
\usepackage[a4paper,margin=1in]{geometry}
\usepackage{microtype}
\usepackage{graphicx}
\usepackage{xcolor}
\usepackage{booktabs}
\usepackage{longtable}
\usepackage{enumitem}
\usepackage{hyperref}
\usepackage{listings}
\usepackage{amsmath}

\hypersetup{
  colorlinks=true,
  linkcolor=blue!60!black,
  urlcolor=blue!60!black,
  citecolor=blue!60!black,
  pdftitle={Multimodal AI Study Companion - Track D},
  pdfauthor={Multimodal AI Hackathon Track D Team}
}

\definecolor{codebg}{RGB}{245,247,250}
\definecolor{codecomment}{RGB}{80,100,80}
\definecolor{codekeyword}{RGB}{0,70,140}

\lstset{
  basicstyle=\ttfamily\small,
  backgroundcolor=\color{codebg},
  frame=single,
  rulecolor=\color{black!15},
  breaklines=true,
  columns=fullflexible,
  showstringspaces=false,
  commentstyle=\color{codecomment},
  keywordstyle=\color{codekeyword}\bfseries
}

\title{\textbf{Multimodal AI Study Companion}\\
\large Hackathon Track D Project Documentation}
\author{Multimodal AI Hackathon Track D Team}
\date{October 2026}

\begin{document}

\maketitle
\tableofcontents
\newpage

\section{Executive Summary}

The Multimodal AI Study Companion is an AI-powered learning workspace for
lecture videos, textbooks, slides, and images. The system converts uploaded
learning material into searchable, source-aware knowledge. Students can ask
questions, receive answers grounded in their own course material, open the
supporting page or video timestamp, and practice with adaptive quizzes.

The current implementation combines a Next.js frontend, a FastAPI backend,
Supabase PostgreSQL and Storage, a Python ingestion pipeline, ChromaDB hybrid
retrieval, and Groq/Gemini language model integrations. Authentication is
performed using a JWT issued by the companion authentication server.

The central design principle is:

\begin{quote}
\emph{Store course knowledge in retrieval, and use the language model for
grounded tutoring behavior.}
\end{quote}

This approach avoids placing private course documents into model weights and
allows uploaded files to be replaced or deleted.

\section{Project Objectives}

\begin{enumerate}[leftmargin=*]
  \item Ingest PDFs, scanned PDFs, images, and lecture videos.
  \item Preserve source metadata such as page numbers and video timestamps.
  \item Support permission-scoped retrieval for personal and community files.
  \item Generate grounded answers with verifiable source citations.
  \item Provide deep links from citations to a PDF page or video timestamp.
  \item Generate course-material-based quizzes with explanations.
  \item Track learner mastery using a client-side Bayesian Knowledge Tracing
        (BKT)-style model.
  \item Support the complete file lifecycle, including removal of indexed
        content after a student deletes an uploaded file.
\end{enumerate}

\section{System Architecture}

\subsection{High-Level Components}

\begin{longtable}{p{0.24\textwidth}p{0.68\textwidth}}
\toprule
\textbf{Component} & \textbf{Responsibility} \\
\midrule
Next.js frontend & Three-pane study workspace, authentication flow,
resource list, chat stream, quiz interface, PDF/video/image viewers, and
mastery visualization. \\
FastAPI backend & Authentication, authorization, resource APIs, chat and quiz
endpoints, signed URLs, background ingestion, and integration with retrieval
and language model services. \\
Supabase PostgreSQL & Communities, memberships, resources, documents, and chat
metadata. \\
Supabase Storage & Private and community file storage using signed URLs. \\
Python ingestion package & File extraction, OCR, transcription, chunking,
embedding, indexing, and indexed-document deletion. \\
ChromaDB & Persistent vector collection containing source-aware document and
transcript chunks. \\
BM25 & Keyword retrieval over permission-filtered chunks. \\
Groq/Gemini & Grounded answer generation and quiz generation. \\
JWT authentication server & Google sign-in integration and signed identity
claims used by the FastAPI backend. \\
\bottomrule
\end{longtable}

\subsection{Request and Data Flow}

\begin{lstlisting}[language={},caption={Upload and question-answering flow}]
Upload:
Browser -> FastAPI upload route
        -> Supabase Storage and metadata row
        -> background ingestion task
        -> extraction/transcription
        -> token-aware chunking
        -> embeddings and ChromaDB upsert

Question:
Browser -> authenticated FastAPI chat route
        -> database membership lookup
        -> permission-scoped vector and BM25 retrieval
        -> reciprocal-rank fusion
        -> Groq/Gemini prompt with source metadata
        -> citation validation
        -> SSE answer stream or JSON response
\end{lstlisting}

\subsection{Repository Layout}

\begin{lstlisting}[language={},caption={Relevant repository directories}]
backend/       FastAPI application, routers, authentication, LLM and bridge
frontend/      Next.js application and study workspace
ingestion/     Extraction, transcription, chunking and ChromaDB indexing
db/            Prisma schema and database migrations
auth-server/   Authentication service
README.md      Project overview and setup information
\end{lstlisting}

\section{Multimodal Ingestion Pipeline}

\subsection{PDF Extraction}

The PDF parser uses PyMuPDF. It first attempts native text extraction for each
page. If a page has no selectable text, the page is rendered to a PNG image at
high resolution and sent to the Gemini OCR helper. The output preserves:

\begin{itemize}[leftmargin=*]
  \item original file name;
  \item one-based page number;
  \item extracted text;
  \item a boolean indicating whether text was found.
\end{itemize}

This two-stage approach supports both normal text PDFs and scanned documents.

\subsection{Image Extraction}

PNG, JPG, and JPEG files are represented as one-page documents. The image bytes
are passed to Gemini OCR with the corresponding MIME type. The resulting
document uses the same page-level shape as PDF extraction, which allows both
formats to share the document chunking and indexing path.

\subsection{Video Transcription}

The video transcriber uses \texttt{faster-whisper}. MP4, MKV, MOV, AVI, and
WebM files are supported by the ingestion script. Each non-empty transcript
segment stores:

\begin{lstlisting}[language=Python]
{
    "start_time": 12.345,
    "end_time": 18.920,
    "text": "The explanation of the concept starts here."
}
\end{lstlisting}

The detected language, language probability, model size, and timestamped
segments are written to a JSON transcript.

\subsection{Chunking}

Document and transcript chunking uses the \texttt{cl100k\_base} tokenizer.
The default chunk size is 400 tokens with 50 tokens of overlap.

For documents, chunks remain associated with their page number. For
transcripts, chunks combine consecutive segments and preserve the first
segment's start time and the last segment's end time. Oversized transcript
segments are split into smaller token-sized pieces.

\begin{longtable}{p{0.28\textwidth}p{0.60\textwidth}}
\toprule
\textbf{Metadata field} & \textbf{Meaning} \\
\midrule
\texttt{file\_name} & Original source file name. \\
\texttt{page\_number} & PDF or image page location; null for video chunks. \\
\texttt{start\_time} & Video start time in seconds; null for document chunks. \\
\texttt{end\_time} & Video end time in seconds; null for document chunks. \\
\texttt{chunk\_index} & Chunk position within the indexed source. \\
\texttt{document\_id} & Stable database/resource identifier. \\
\texttt{owner\_id} & Uploading user's identity. \\
\texttt{access\_scope} & Either \texttt{private} or \texttt{community}. \\
\texttt{community\_id} & Community identifier when sharing is enabled. \\
\bottomrule
\end{longtable}

\section{ChromaDB Indexing and Retrieval}

\subsection{Embedding and Indexing}

The ingestion package uses the \texttt{all-MiniLM-L6-v2} sentence-transformer
to embed non-empty chunks. Embeddings are normalized and stored in the
persistent Chroma collection \texttt{course\_material\_scoped}.

Each Chroma record has a stable identifier derived from the database document
identifier:

\begin{lstlisting}[language=Python]
chunk_id = f"{document_id}_{chunk_index}_{page_number}_{start_time}"
\end{lstlisting}

The \texttt{document\_id} is also stored explicitly in metadata. The metadata
field is the authoritative key for deleting all chunks belonging to a source.

\subsection{Permission-Scoped Hybrid Retrieval}

The retrieval process is:

\begin{enumerate}[leftmargin=*]
  \item Build a private-document filter for the authenticated user.
  \item Query the database for communities the user belongs to.
  \item Add authorized community documents to the access filter.
  \item Run Chroma vector retrieval using the combined filter.
  \item Run BM25 keyword retrieval only over the same filtered records.
  \item Combine vector and keyword rankings using Reciprocal Rank Fusion.
  \item Convert the selected records into backend source metadata.
\end{enumerate}

This prevents unauthorized private or community chunks from entering either
the semantic or keyword search path.

\subsection{Citation Handling}

The LLM prompt includes a required source tag for every retrieved chunk. The
citation guardrail scans generated answers, removes unmatched citation tags,
and returns validated citation objects. PDF citations contain page numbers;
video citations contain timestamp ranges.

The frontend resolves citations to a source and opens the appropriate viewer.
Video citations seek the Video.js player, while PDF citations navigate to the
referenced page. The current implementation supports the viewer's highlight
metadata contract, although automatic OCR/text bounding-box generation is not
yet complete.

\section{Storage, Metadata, and Access Control}

\subsection{Supabase Storage}

The application uses separate buckets:

\begin{itemize}[leftmargin=*]
  \item \texttt{private-materials} for personal files;
  \item \texttt{community-materials} for community files.
\end{itemize}

Storage paths include an identity and a stable resource identifier. The
backend returns short-lived signed URLs only after checking ownership or
community membership.

\subsection{Database Model}

The database contains the following principal entities:

\begin{itemize}[leftmargin=*]
  \item \textbf{Community}: shared learning space and invite code;
  \item \textbf{CommunityMember}: user membership relation;
  \item \textbf{Resource}: community resource metadata;
  \item \textbf{Document}: personal or indexed resource metadata;
  \item \textbf{Chat}: chat-session metadata.
\end{itemize}

The current processing states are:

\begin{center}
\texttt{uploaded} $\rightarrow$ \texttt{processing} $\rightarrow$
\texttt{completed}
\end{center}

Failures use \texttt{failed}. Deletion uses \texttt{deleting} followed by
\texttt{deleted}.

\section{Deletion and Index Consistency}

\subsection{Why Deletion Needs Coordination}

Ingestion runs as a background task. Without coordination, the following race
could occur:

\begin{enumerate}[leftmargin=*]
  \item A file starts OCR or transcription.
  \item The student deletes the file.
  \item The delete route removes existing Chroma chunks.
  \item The old background task finishes and inserts the chunks again.
\end{enumerate}

\subsection{Implemented Deletion Flow}

The current worktree contains the following deletion flow:

\begin{enumerate}[leftmargin=*]
  \item The backend verifies ownership or community membership.
  \item The metadata record is marked \texttt{deleting}.
  \item All Chroma records matching the stable \texttt{document\_id} are
        removed.
  \item The Supabase Storage object is removed.
  \item The metadata record is marked \texttt{deleted}.
  \item Deleted resources are excluded from normal resource listings.
\end{enumerate}

The background ingestion bridge checks whether the source record is missing,
\texttt{deleting}, or \texttt{deleted}. A final cancellation check occurs after
extraction/chunking and immediately before Chroma insertion.

The API routes are:

\begin{lstlisting}[language={}]
DELETE /api/personal/resources/{resource_id}
DELETE /api/communities/{community_id}/resources/{resource_id}
\end{lstlisting}

This design preserves document isolation and avoids deleting chunks by
non-unique file names.

\section{Backend Application}

The FastAPI application exposes:

\begin{itemize}[leftmargin=*]
  \item health check endpoint;
  \item synchronous grounded chat endpoint;
  \item server-sent event chat endpoint;
  \item quiz generation endpoint;
  \item personal resource upload, listing, signed URL, and deletion;
  \item community creation, joining, membership, resource upload, signed URL,
        and deletion;
  \item chat-session listing and deletion.
\end{itemize}

The backend authenticates requests using a bearer JWT. The service-role
Supabase client is kept on the backend, while the frontend receives only
short-lived signed resource URLs.

\section{Frontend Application}

\subsection{Study Workspace}

The frontend provides a three-pane workflow:

\begin{enumerate}[leftmargin=*]
  \item \textbf{Sources}: personal and community materials;
  \item \textbf{Chat}: streaming grounded questions and answers;
  \item \textbf{Viewer}: PDF, video, and image source viewer.
\end{enumerate}

\subsection{Viewers and Deep Links}

The PDF viewer uses \texttt{react-pdf} and supports page navigation. The
Video.js player supports programmatic seeking to a citation timestamp. Images
are displayed directly from signed Storage URLs.

When a selected resource is deleted, the viewer context detects that the
resource is no longer present and clears the open viewer target.

\subsection{Assessment and Mastery}

The quiz interface supports multiple-choice questions, answer feedback,
explanations, source-citation text, and mastery updates. The client-side BKT
helpers classify topic mastery and adapt quiz length.

The current implementation provides the user-interface and local learner
model foundation. Formal simulated-student results and a backend-persisted
learner profile remain future evaluation work.

\section{Current Implementation Status}

\begin{longtable}{p{0.34\textwidth}p{0.20\textwidth}p{0.38\textwidth}}
\toprule
\textbf{Feature} & \textbf{Status} & \textbf{Notes} \\
\midrule
Native PDF extraction & Implemented & PyMuPDF page extraction. \\
Scanned PDF OCR & Implemented & Gemini OCR fallback. \\
Image OCR & Implemented & PNG/JPG/JPEG support. \\
Video transcription & Implemented & Timestamped faster-whisper output. \\
Token-aware chunking & Implemented & Documents and transcripts. \\
ChromaDB indexing & Implemented & Persistent embeddings and metadata. \\
Hybrid retrieval & Implemented & Vector search, BM25, and RRF. \\
Permission filtering & Implemented & Private owner and authorized communities. \\
Supabase Storage & Implemented & Uploads and signed URLs. \\
Personal deletion & Implemented in current worktree & Chroma, Storage, and metadata cleanup. \\
Community deletion & Implemented in current worktree & Membership-protected cleanup. \\
Deletion race protection & Implemented in current worktree & Processing state and cancellation checks. \\
Citation validation & Partially implemented & Non-streaming validation is present. \\
Streaming citation events & Incomplete & Frontend/backend event behavior needs alignment. \\
Grounded quiz retrieval & Incomplete & Quiz route needs topic-based server retrieval. \\
BKT personalization & Prototype & Client-side state and adaptive length. \\
Formal evaluation report & Not yet implemented & Requires benchmark data and metrics. \\
\bottomrule
\end{longtable}

\section{Testing and Verification}

The repository includes tests for:

\begin{itemize}[leftmargin=*]
  \item PDF extraction and OCR fallback;
  \item image parsing;
  \item document routing;
  \item document and transcript chunking;
  \item ingestion pipeline routing;
  \item transcript metadata preservation;
  \item access-control-related behavior;
  \item storage and retrieval helper behavior.
\end{itemize}

The changed Python files were syntax-checked with
\texttt{python -m compileall}. Whitespace validation passed with
\texttt{git diff --check}. In the development environment used for this
documentation, \texttt{pytest} was not installed and the frontend
\texttt{node\_modules} directory was unavailable, so the complete test suite,
frontend lint, and frontend production build could not be executed there.

Before deployment, the database migration
\texttt{db/prisma/migrations/0004\_resource\_processing\_status/migration.sql}
must be applied to the Supabase database.

\subsection{Ingestion System Evaluation}

The repository now includes a reproducible retrieval evaluation harness under
\texttt{ingestion/evaluation}. The labeled fixture set covers:

\begin{itemize}[leftmargin=*]
  \item a page-cited PDF question;
  \item a timestamp-cited video question;
  \item an authorized community-resource question;
  \item a private-document isolation case.
\end{itemize}

Run the evaluation from the repository root with:

\begin{lstlisting}[language=bash]
python -m ingestion.evaluation.evaluate_retrieval --top-k 5
\end{lstlisting}

The runner temporarily indexes the fixture corpus, reports Recall@k, mean
reciprocal rank, source-metadata accuracy, unsupported-query context count,
and unauthorized-retrieval count, and removes the fixture records in a
\texttt{finally} block. It therefore does not intentionally leave benchmark
records in the persistent Chroma collection. The fixture set is a baseline
evaluation and should be expanded with course-specific and multimodal cases
before reporting final system results.

\section{Known Limitations}

\begin{enumerate}[leftmargin=*]
  \item Streaming chat currently emits token events and completion events, but
        citation-event delivery is not yet fully aligned with the frontend
        citation contract.
  \item The quiz frontend can request questions without first performing
        topic-based permission-scoped retrieval.
  \item Community discovery is partly remembered in browser local storage
        rather than being fully reconstructed from a backend membership list.
  \item PDF and image OCR text is preserved, but exact visual text coordinates
        are not generated by the ingestion pipeline.
  \item Retrieval failures should be distinguished more clearly from an empty
        retrieval result.
  \item Standalone ingestion storage helpers and backend upload routes contain
        overlapping storage logic and should eventually be consolidated.
  \item Formal measurements for grounding, assessment quality, personalization,
        latency, and deletion consistency are not yet reported.
\end{enumerate}

\section{Recommended Evaluation Plan}

To evaluate the system rigorously, create a labeled test set with the
following fields:

\begin{lstlisting}[language=json,caption={Suggested evaluation record}]
{
  "id": "video-001",
  "question": "What is the purpose of the learning rate?",
  "topic": "gradient descent",
  "answerable": true,
  "expected_sources": [
    {"file": "lecture.mp4", "start_time": "02:15"}
  ],
  "expected_concepts": ["step size", "parameter update"]
}
\end{lstlisting}

Recommended metrics are:

\begin{itemize}[leftmargin=*]
  \item Retrieval Recall@1, Recall@3, Recall@5, and MRR;
  \item citation precision and citation recall;
  \item grounded-answer faithfulness;
  \item unsupported-query refusal rate;
  \item quiz correctness, source support, topic relevance, and duplicate rate;
  \item learner accuracy and mastery change across simulated quiz runs;
  \item upload-to-index latency, query latency, and failure rate;
  \item deletion consistency across database, Storage, and ChromaDB.
\end{itemize}

Results should include a baseline, the improved system, examples of failures,
and known limitations. A small, reproducible test set is preferable to an
unverified accuracy claim.

\section{Recommended Next Steps}

\begin{enumerate}[leftmargin=*]
  \item Apply and verify the resource-processing database migration.
  \item Add integration tests for upload, indexing, deletion, and re-query.
  \item Complete server-side topic retrieval for grounded quizzes.
  \item Align streaming citation events with the frontend contract.
  \item Add explicit unsupported-query refusal tests.
  \item Run simulated learner experiments and report mastery gains.
  \item Add a backend endpoint for listing all user communities.
  \item Consolidate duplicate Storage/upload helper logic.
  \item Run frontend lint/build after installing the frontend dependencies.
  \item Record a demo showing upload, citation navigation, adaptive practice,
        unsupported-query refusal, and deletion from the knowledge base.
\end{enumerate}

\section{Conclusion}

The current project has a strong multimodal RAG foundation. It already
supports extraction from documents, images, and videos; source-aware
chunking; hybrid permission-scoped retrieval; Supabase Storage; citation
navigation; and adaptive quiz interfaces. The newly added deletion lifecycle
keeps the persistent knowledge index consistent when an uploaded resource is
removed.

The highest-value remaining work is not model fine-tuning. It is completing
grounded quiz retrieval, making streaming citations fully reliable, and
producing a rigorous evaluation report with measurable grounding and learner
improvement results.

\end{document}
